\section{Methods}

\subsection{Data and response}

A document is one movie's complete overview. Of the 45,466 source rows, we keep
English-original-language films with valid movie IDs and usable overviews.
We normalize whitespace, remove known placeholders, collapse equivalent
records, and exclude conflicting movie IDs. This leaves 32,105 films for topic
analysis. We treat zero budget, runtime, and revenue as missing. There are
6,348 cleaned films with positive recorded revenue and a valid release date.
The predictive comparison also requires a listed production company and
release-year CPI coverage. Both revenue targets use the same films.

The primary response is the natural log of recorded revenue. The source
contains no per-film currency or geographic gross definition. In a sensitivity
analysis, we assume comparable nominal USD and multiply recorded revenue by the
ratio of 2016 CPI-U to release-year CPI-U before taking logs. Annual CPI-U is
the mean of the twelve monthly observations from the BLS series obtained
through FRED \cite{cpi}. Release year approximates the timing of receipts.

\subsection{Text representations and topic analysis}

We form a TF--IDF matrix from lowercase alphabetic unigrams of at least two
letters. We remove English stop words and keep terms appearing in at least two
documents. The descriptive corpus has 30,054 retained terms. TF--IDF preserves
specific words that may distinguish movies within the same genre.

Nonnegative matrix factorization (NMF) approximates the TF--IDF matrix as a
product of movie--topic and topic--word matrices \cite{lee1999}. The five-topic
specification keeps the topic table readable while allowing each overview to
load on several themes. For descriptive assignments, we scale the movie
weights by the corresponding topic-vector norms and use the largest weight.
We select each representative film by cosine similarity to its dominant
topic. We refit twice with perturbed initializations, match topics by cosine
similarity, and compare their ten highest-weight words.

The frozen \texttt{all-MiniLM-L6-v2} sentence encoder supplies a
384-dimensional embedding \cite{minilm,reimers2019}.
The model uses its own tokenizer on the intact overview. For overviews above
its context limit, we encode token windows, pool their vectors by token count,
and normalize. This representation tests whether contextual information adds
signal beyond the lexical features and compressed topic weights.

\subsection{Controls and chronological evaluation}

The controls are release year and month, runtime, multi-label genre, collection
membership as a franchise proxy, and multi-label production-company indicators.
Year and runtime are standardized, missing runtime is imputed with the training
median, and month is one-hot encoded. Each company observed on at least five
training films gets its own indicator. Other companies share a rare-company
indicator. We use ridge regression with a fixed penalty of 1 for all
comparisons. We add each plot representation separately and omit one control
block at a time in the controls-only model to measure its contribution given
the remaining blocks.

We evaluate releases in 2005--2008, 2009--2011, 2012--2013, and 2014--2016,
using only earlier films for training. The windows contain 806, 678, 461, and
715 films, respectively. The final window was inspected during model
development. We fit TF--IDF, predictive NMF, and metadata transforms within
each training fold. The descriptive NMF fit is separate. We exclude training
overviews that exactly match a test overview after whitespace and case
normalization.

For film $i$, let $e_i^{(0)}$ and $e_i^{(1)}$ be the absolute log-revenue
errors for the controls and comparison models. We report the paired mean
$\Delta\mathrm{MAE} = n^{-1}\sum_i (e_i^{(1)}-e_i^{(0)})$.
We form pointwise 95\% percentile intervals from 2,000 paired bootstrap
resamples of films within release year. These intervals are conditional on
the observed years and fitted models. We display genre and company-support
slices with at least 30 films and interpret them descriptively.
